\documentclass[a4paper]{arrowhead}

\usepackage[yyyymmdd]{datetime}
\usepackage{etoolbox}
\usepackage[utf8]{inputenc}
\usepackage[T1]{fontenc}
\usepackage{multirow}
\usepackage{longtable}
\usepackage{tikz}
\usepackage{adjustbox}
\AtBeginEnvironment{longtable}{\small}
\BeforeBeginEnvironment{tikzpicture}{\begin{adjustbox}{max width=\textwidth}}
\AfterEndEnvironment{tikzpicture}{\end{adjustbox}}
\usetikzlibrary{positioning,shapes.multipart,arrows.meta}

\renewcommand{\dateseparator}{-}
\setlength{\emergencystretch}{3em}
\lstset{language=HTTP,breakatwhitespace=false,columns=fullflexible}

\newcommand{\fparam}[1]{\textit{\textcolor{ArrowheadBlue}{#1}}}

%% Special references
\newcommand{\mref}[2][]{{\textcolor{ArrowheadPurple}{\hyperref[sec:model:#2]{\ifstrempty{#1}{#2}{#1}}}}}
\newcommand{\oref}[1]{{\textcolor{ArrowheadBlue}{\hyperref[sec:operations:#1]{#1}}}}
\newcommand{\pdef}[1]{{\textcolor{ArrowheadGrey}{#1 \label{sec:model:primitives:#1} \label{sec:model:primitives:#1s}}}}
\newcommand{\pref}[1]{{\textcolor{ArrowheadGrey}{\hyperref[sec:model:primitives:#1]{#1}}}}

% as template_idd.tex; the input and output arguments are passed already formatted (\mref or \pref)
\newrobustcmd\osubsection[5]{
  \addtocounter{subsection}{1}
  \addcontentsline{toc}{subsection}{\protect\numberline{\thesubsection}operation \textcolor{ArrowheadBlue}{#1}}
  \subsection*{
    \thesubsection\quad
    #2 \textcolor{ArrowheadPurple}{#3} \\
    \small
    \hspace*{0.075\textwidth}\begin{minipage}{0.1\textwidth}
      \vspace*{1mm}
      Operation: \\
      \notblank{#4}{Input: \\}{}
      \notblank{#5}{Output: \\}{}
    \end{minipage}%
    \begin{minipage}{0.825\textwidth}
      \vspace*{1mm}
      \textcolor{ArrowheadBlue}{#1} \\
      \notblank{#4}{#4 \\}{}
      \notblank{#5}{#5 \\}{}
    \end{minipage}
  }
  \label{sec:operations:#1}
}
\newrobustcmd\msubsection[2]{
  \addtocounter{subsection}{1}
  \addcontentsline{toc}{subsection}{\protect\numberline{\thesubsection}#1 \textcolor{ArrowheadPurple}{#2}}
  \subsection*{\thesubsection\quad#1 \textcolor{ArrowheadPurple}{#2}}
  \label{sec:model:#2}
}
%%

\begin{document}
\sloppy

%% Arrowhead Document Properties
\ArrowheadTitle{authorization HTTP/TLS/JSON}
\ArrowheadServiceID{authorization}
\ArrowheadType{Interface Design Description}
\ArrowheadTypeShort{IDD}
\ArrowheadVersion{5.2.2}
\ArrowheadDate{2026-10-22}
\ArrowheadAuthor{Jerker Delsing}
\ArrowheadStatus{RELEASE CANDIDATE}
\ArrowheadContact{jerker.delsing@ltu.se}
\ArrowheadFooter{\href{www.arrowhead.eu}{www.arrowhead.eu}}
\ArrowheadSetup
%%

%% Front Page
\begin{center}
  \vspace*{1cm}
  \huge{\arrowtitle}

  \vspace*{0.2cm}
  \LARGE{\arrowtype}
  \vspace*{1cm}

  \Large{Service ID: \textit{"\arrowid"}}

  \vspace*{\fill}

  % Front Page Image

  \vspace*{1cm}
  \vspace*{\fill}

  % Front Page Abstract
  \begin{abstract}
    This document describes the HTTP/1.1 protocol with TLS payload security and JSON payload encoding variant of the \texttt{authorization} service of the ConsumerAuthorization system.
  \end{abstract}

  \vspace*{1cm}
\end{center}

\newpage
%%

%% Table of Contents
\tableofcontents
\newpage
%%

\section{Overview}
\label{sec:overview}

This document describes the \texttt{authorization} service interface of the ConsumerAuthorization system, which provides provider-level authorization policies: a provider declares which consumers may use its service (or event type), and any system can ask whether a consumer is authorized.

It's implemented using protocol, encoding as stated in below table:
\begin{table}[ht!]
  \centering
  \begin{tabular}{|l|l|l|}
    \hline
    \rowcolor{gray!33} Profile type & Type & Version \\ \hline
    Transfer protocol & HTTP & 1.1 \cite{rfc9112} \\ \hline
    Data encryption & TLS & 1.3 \cite{rfc8446}, 1.2 \cite{rfc5246} \\ \hline
    Encoding & JSON & RFC 8259 \cite{rfc8259} \\ \hline
    Compression & N/A & - \\ \hline
    Semantics & N/A & - \\ \hline
    Ontology & N/A & - \\ \hline
  \end{tabular}
  \caption{Communication and semantics details used for the \texttt{authorization} service interface}
  \label{tab:comunication_semantics_profile}
\end{table}

This document provides the Interface Design Description IDD to the \textit{authorization -- Service Description} document \cite{ca-authorization-sd}.
For further details about how this service is meant to be used, please consult that document.

The rest of this document describes how to realize the \texttt{authorization} service
interface in detail. Both in terms of its operations (Section \ref{sec:operations}) and its information model (Section \ref{sec:model}).

\subsection{Interface registration}

The service is registered in the ServiceRegistry with the interface template \texttt{generic\_\allowbreak{}https} (\texttt{generic\_\allowbreak{}http} when TLS is disabled), access port 8445 and base path \texttt{/\allowbreak{}consumerauthorization/\allowbreak{}authorization}. The \texttt{operations} property of the interface lists the operations below with their method and path relative to the base path (path parameters are placeholders such as \texttt{\{id\}} that the consumer fills in):

\noindent\begin{longtable}{| p{5.2cm} | p{1.6cm} | p{8.4cm} |} \hline
\rowcolor{gray!33} Operation & Method & Path \\ \hline \endhead
\texttt{grant} & POST & \texttt{/\allowbreak{}grant} \\ \hline
\texttt{revoke} & DELETE & \texttt{/\allowbreak{}revoke/\allowbreak{}\{instance\allowbreak{}Id\}} \\ \hline
\texttt{lookup} & POST & \texttt{/\allowbreak{}lookup} \\ \hline
\texttt{verify} & POST & \texttt{/\allowbreak{}verify} \\ \hline
\end{longtable}

\newpage

\section{Service Operations}
\label{sec:operations}

The interfaces of the \texttt{authorization} service, its operations, data models and implementation are provided below.
In particular, each subsection first names the HTTP method and path used to call the operation, after which it names the abstract operation from the \texttt{authorization} SD document, as well as input and output types.

All operations require the requester's identity, presented according to the provider's authentication policy: the client certificate (mutual TLS) and, with the \texttt{outsourced} policy, \texttt{Authorization: Bearer IDENTITY-TOKEN//<token>}; systems holding the Authentication secret key send \texttt{Bearer AUTHENTICATOR-KEY//<system>//<hash>}, and with the \texttt{declared} policy \texttt{Bearer SYSTEM//<system>} is sent. Request and response bodies are \texttt{application/json}.
Errors are answered with an \mref{ErrorMessage} body and the status codes \texttt{400 Bad Request} (invalid parameter), \texttt{401 Unauthorized} (missing or invalid identity), \texttt{403 Forbidden} (not permitted, e.g. by the management or blacklist filter), \texttt{404 Not Found}, \texttt{408 Request Timeout}, \texttt{423 Locked}, \texttt{500 Internal Server Error} and \texttt{503 Service Unavailable} (a system it depends on failed).

\osubsection{grant}{POST}{/\allowbreak{}consumerauthorization/\allowbreak{}authorization/\allowbreak{}grant}{\mref[Authorization\allowbreak{}Grant\allowbreak{}Request]{AuthorizationGrantRequest}}{\mref[Authorization\allowbreak{}Policy\allowbreak{}Response]{AuthorizationPolicyResponse}}

Creates or replaces the provider-level policy of the requesting provider for one target. Returns 201 for a new policy, 200 if an identical one exists.

A successful call returns \texttt{201 Created} for a new entry, \texttt{200 OK} if an identical entry already existed with the output in the body.

\begin{lstlisting}[label={lst:grant},caption={An \oref{grant} invocation.}]
POST /consumerauthorization/authorization/grant HTTP/1.1
Host: consumerauthorization:8445
Authorization: Bearer IDENTITY-TOKEN//<token>
Content-Type: application/json

{
  "cloud": "LOCAL",
  "targetType": "SERVICE_DEF",
  "target": "temperatureReading",
  "description": "Displays may read the temperature",
  "defaultPolicy": {
    "policyType": "WHITELIST",
    "policyList": [
      "TemperatureDisplay"
    ],
    "policyMetadataRequirement": {
      "unit": "celsius"
    }
  },
  "scopedPolicies": {
    "read": {
      "policyType": "ALL"
    }
  }
}
\end{lstlisting}

\begin{lstlisting}[label={lst:grant-resp},caption={An \oref{grant} response.}]
{
  "instanceId": "TemperatureSensor|temperatureReading|1.0.0",
  "level": "MGMT",
  "cloud": "LOCAL",
  "provider": "TemperatureSensor",
  "targetType": "SERVICE_DEF",
  "target": "temperatureReading",
  "description": "Display may read",
  "defaultPolicy": {
    "policyType": "ALL",
    "policyList": [
      "TemperatureDisplay"
    ],
    "policyMetadataRequirement": {
      "unit": "celsius"
    }
  },
  "scopedPolicies": {
    "read": {
      "policyType": "ALL"
    }
  },
  "createdBy": "Sysop",
  "createdAt": "2026-10-22T12:00:00Z"
}
\end{lstlisting}

\osubsection{revoke}{DELETE}{/\allowbreak{}consumerauthorization/\allowbreak{}authorization/\allowbreak{}revoke/\allowbreak{}\fparam{\{instanceId\}}}{\pref{String}}{}

Removes a provider-level policy of the requester by instance identifier. Returns 200 if removed, 204 if it did not exist.

Path parameter: \fparam{instanceId}.

A successful call returns \texttt{200 OK} if the entry was removed, \texttt{204 No Content} if there was nothing to remove.

\begin{lstlisting}[label={lst:revoke},caption={An \oref{revoke} invocation.}]
DELETE /consumerauthorization/authorization/revoke/<instanceId> HTTP/1.1
Host: consumerauthorization:8445
Authorization: Bearer IDENTITY-TOKEN//<token>
\end{lstlisting}

\osubsection{lookup}{POST}{/\allowbreak{}consumerauthorization/\allowbreak{}authorization/\allowbreak{}lookup}{\mref[Authorization\allowbreak{}Lookup\allowbreak{}Request]{AuthorizationLookupRequest}}{\mref[Authorization\allowbreak{}Policy\allowbreak{}List\allowbreak{}Response]{AuthorizationPolicyListResponse}}

Returns the requester's provider-level policies matching the filter.

A successful call returns \texttt{200 OK} with the output in the body.

\begin{lstlisting}[label={lst:lookup},caption={An \oref{lookup} invocation.}]
POST /consumerauthorization/authorization/lookup HTTP/1.1
Host: consumerauthorization:8445
Authorization: Bearer IDENTITY-TOKEN//<token>
Content-Type: application/json

{
  "instanceIds": [
    "TemperatureSensor|temperatureReading|1.0.0"
  ],
  "cloudIdentifiers": [
    "string"
  ],
  "targetNames": [
    "string"
  ],
  "targetType": "SERVICE_DEF"
}
\end{lstlisting}

\begin{lstlisting}[label={lst:lookup-resp},caption={An \oref{lookup} response.}]
{
  "entries": [
    {
      "instanceId": "TemperatureSensor|temperatureReading|1.0.0",
      "level": "MGMT",
      "cloud": "LOCAL",
      "provider": "TemperatureSensor",
      "targetType": "SERVICE_DEF",
      "target": "temperatureReading",
      "description": "Display may read",
      "defaultPolicy": {
        "policyType": "ALL",
        "policyList": [
          "TemperatureDisplay"
        ],
        "policyMetadataRequirement": {
          "unit": "celsius"
        }
      },
      "scopedPolicies": {
        "read": {
          "policyType": "ALL"
        }
      },
      "createdBy": "Sysop",
      "createdAt": "2026-10-22T12:00:00Z"
    }
  ],
  "count": 1
}
\end{lstlisting}

\osubsection{verify}{POST}{/\allowbreak{}consumerauthorization/\allowbreak{}authorization/\allowbreak{}verify}{\mref[Authorization\allowbreak{}Verify\allowbreak{}Request]{AuthorizationVerifyRequest}}{\pref{Boolean}}

Evaluates whether the consumer may use the target (and operation) of the provider: management-level policies first, provider-level policies if none applies; no applicable policy means denied. Returns a boolean.

A successful call returns \texttt{200 OK} with the output in the body.

\begin{lstlisting}[label={lst:verify},caption={An \oref{verify} invocation.}]
POST /consumerauthorization/authorization/verify HTTP/1.1
Host: consumerauthorization:8445
Authorization: Bearer IDENTITY-TOKEN//<token>
Content-Type: application/json

{
  "provider": "TemperatureSensor",
  "consumer": "TemperatureDisplay",
  "cloud": "LOCAL",
  "targetType": "SERVICE_DEF",
  "target": "temperatureReading",
  "scope": "read"
}
\end{lstlisting}

\begin{lstlisting}[label={lst:verify-resp},caption={An \oref{verify} response.}]
true
\end{lstlisting}

\newpage

\section{Data Models}
\label{sec:model}

Here, all data objects that can be part of the service calls associated with this service are listed in alphabetic order.
Note that each subsection, which describes one type of object, begins with the \textit{struct} keyword, which is meant to denote a JSON \pref{Object} that must contain certain fields, or names, with values conforming to explicitly named types.
Fields that are not needed may be omitted or \texttt{null} in requests; empty fields are omitted in responses.
As a complement to the primary types defined in this section, there is also a list of secondary types in Section \ref{sec:model:primitives}, which are used to represent things like hashes, identifiers and texts.

\msubsection{struct}{AuthorizationGrantRequest}

A provider-level authorization policy that the provider sets for its own service.


\noindent\begin{tabularx}{\textwidth}{| >{\raggedright\arraybackslash}p{3.9cm} | >{\raggedright\arraybackslash\small}p{5.3cm} | X |} \hline
\rowcolor{gray!33} Object Field & Value Type & Description \\ \hline
"cloud" & \pref{String} & Cloud identifier of the consumer cloud the policy applies to; \texttt{LOCAL} for the own cloud. \\ \hline
"target\allowbreak{}Type" & \pref{String} & Target type: \texttt{SERVICE\_\allowbreak{}DEF} or \texttt{EVENT\_\allowbreak{}TYPE}. \\ \hline
"target" & \pref{String} & Service definition (or event type) the policy protects. \\ \hline
"description" & \pref{String} & Free-text description. \\ \hline
"default\allowbreak{}Policy" & \mref[Authorization\allowbreak{}Policy\allowbreak{}Request]{AuthorizationPolicyRequest} & Policy that applies to all operations without a scoped policy. \\ \hline
"scoped\allowbreak{}Policies" & \pref{Object}$<$\allowbreak{}\mref[Authorization\allowbreak{}Policy\allowbreak{}Request]{AuthorizationPolicyRequest}$>$ & Policies for individual operations, keyed by operation name. \\ \hline
\end{tabularx}



\msubsection{struct}{AuthorizationLookupRequest}

Filter for the requester's own provider-level policies.


\noindent\begin{tabularx}{\textwidth}{| >{\raggedright\arraybackslash}p{3.9cm} | >{\raggedright\arraybackslash\small}p{5.3cm} | X |} \hline
\rowcolor{gray!33} Object Field & Value Type & Description \\ \hline
"instance\allowbreak{}Ids" & \pref{Array}$<$\allowbreak{}\pref{String}$>$ & Policy instance identifiers to match. \\ \hline
"cloud\allowbreak{}Identifiers" & \pref{Array}$<$\allowbreak{}\pref{String}$>$ & Cloud identifiers to match. \\ \hline
"target\allowbreak{}Names" & \pref{Array}$<$\allowbreak{}\pref{String}$>$ & Target names to match. \\ \hline
"target\allowbreak{}Type" & \pref{String} & Target type: \texttt{SERVICE\_\allowbreak{}DEF} or \texttt{EVENT\_\allowbreak{}TYPE}. \\ \hline
\end{tabularx}



\msubsection{struct}{AuthorizationPolicy}

One policy rule.


\noindent\begin{tabularx}{\textwidth}{| >{\raggedright\arraybackslash}p{3.9cm} | >{\raggedright\arraybackslash\small}p{5.3cm} | X |} \hline
\rowcolor{gray!33} Object Field & Value Type & Description \\ \hline
"policy\allowbreak{}Type" & \pref{AuthorizationPolicyType} & Policy type: \texttt{ALL}, \texttt{WHITELIST}, \texttt{BLACKLIST} or \texttt{SYS\_\allowbreak{}METADATA}. \\ \hline
"policy\allowbreak{}List" & \pref{Array}$<$\allowbreak{}\pref{String}$>$ & System names of the white- or blacklist. \\ \hline
"policy\allowbreak{}Metadata\allowbreak{}Requirement" & \mref{MetadataRequirement} & Metadata the consumer system must have (\texttt{SYS\_\allowbreak{}METADATA}). \\ \hline
\end{tabularx}



\msubsection{struct}{AuthorizationPolicyListResponse}

List of authorization policies.


\noindent\begin{tabularx}{\textwidth}{| >{\raggedright\arraybackslash}p{3.9cm} | >{\raggedright\arraybackslash\small}p{5.3cm} | X |} \hline
\rowcolor{gray!33} Object Field & Value Type & Description \\ \hline
"entries" & \pref{Array}$<$\allowbreak{}\mref[Authorization\allowbreak{}Policy\allowbreak{}Response]{AuthorizationPolicyResponse}$>$ & Matching records of the current page. \\ \hline
"count" & \pref{Number} & Total number of matching records (all pages). \\ \hline
\end{tabularx}



\msubsection{struct}{AuthorizationPolicyRequest}

One policy rule.


\noindent\begin{tabularx}{\textwidth}{| >{\raggedright\arraybackslash}p{3.9cm} | >{\raggedright\arraybackslash\small}p{5.3cm} | X |} \hline
\rowcolor{gray!33} Object Field & Value Type & Description \\ \hline
"policy\allowbreak{}Type" & \pref{String} & Policy type: \texttt{ALL}, \texttt{WHITELIST}, \texttt{BLACKLIST} or \texttt{SYS\_\allowbreak{}METADATA}. \\ \hline
"policy\allowbreak{}List" & \pref{Array}$<$\allowbreak{}\pref{String}$>$ & System names of the white- or blacklist. \\ \hline
"policy\allowbreak{}Metadata\allowbreak{}Requirement" & \mref{MetadataRequirement} & Metadata the consumer system must have (\texttt{SYS\_\allowbreak{}METADATA}). \\ \hline
\end{tabularx}



\msubsection{struct}{AuthorizationPolicyResponse}

A stored authorization policy.


\noindent\begin{tabularx}{\textwidth}{| >{\raggedright\arraybackslash}p{3.9cm} | >{\raggedright\arraybackslash\small}p{5.3cm} | X |} \hline
\rowcolor{gray!33} Object Field & Value Type & Description \\ \hline
"instance\allowbreak{}Id" & \pref{String} & Policy instance identifier \texttt{\textless{}MGMT or PR\textgreater{}\textbar{}\allowbreak{}\textless{}cloud\textgreater{}\textbar{}\allowbreak{}\textless{}provider\textgreater{}\textbar{}\allowbreak{}\textless{}target\allowbreak{}Type\textgreater{}\textbar{}\allowbreak{}\textless{}target\textgreater{}}. \\ \hline
"level" & \pref{AuthorizationLevel} & Policy level: \texttt{MGMT} (set by an operator) or \texttt{PROVIDER} (set by the provider itself). \\ \hline
"cloud" & \pref{String} & Cloud identifier of the consumer cloud the policy applies to; \texttt{LOCAL} for the own cloud. \\ \hline
"provider" & \pref{String} & Provider system the policy protects. \\ \hline
"target\allowbreak{}Type" & \pref{AuthorizationTargetType} & Target type: \texttt{SERVICE\_\allowbreak{}DEF} or \texttt{EVENT\_\allowbreak{}TYPE}. \\ \hline
"target" & \pref{String} & Target of the policy or token: a service definition or event type. \\ \hline
"description" & \pref{String} & Free-text description. \\ \hline
"default\allowbreak{}Policy" & \mref{AuthorizationPolicy} & Policy that applies to all operations without a scoped policy. \\ \hline
"scoped\allowbreak{}Policies" & \pref{Object}$<$\allowbreak{}\mref{AuthorizationPolicy}$>$ & Policies for individual operations, keyed by operation name. \\ \hline
"created\allowbreak{}By" & \pref{String} & Name of the system that created the record. \\ \hline
"created\allowbreak{}At" & \pref{DateTime} & UTC time of creation. \\ \hline
\end{tabularx}



\msubsection{struct}{AuthorizationVerifyRequest}

One authorization question: may the consumer use the target (operation) of the provider?


\noindent\begin{tabularx}{\textwidth}{| >{\raggedright\arraybackslash}p{3.9cm} | >{\raggedright\arraybackslash\small}p{5.3cm} | X |} \hline
\rowcolor{gray!33} Object Field & Value Type & Description \\ \hline
"provider" & \pref{String} & Provider system; for \texttt{verify} defaults to the requester. \\ \hline
"consumer" & \pref{String} & Consumer system; for \texttt{verify} defaults to the requester. \\ \hline
"cloud" & \pref{String} & Cloud of the consumer, \texttt{LOCAL} by default. \\ \hline
"target\allowbreak{}Type" & \pref{String} & Target type: \texttt{SERVICE\_\allowbreak{}DEF} or \texttt{EVENT\_\allowbreak{}TYPE}. \\ \hline
"target" & \pref{String} & Target of the policy or token: a service definition or event type. \\ \hline
"scope" & \pref{String} & Operation to check; empty checks the default policy. \\ \hline
\end{tabularx}



\msubsection{struct}{ErrorMessage}




\noindent\begin{tabularx}{\textwidth}{| >{\raggedright\arraybackslash}p{3.9cm} | >{\raggedright\arraybackslash\small}p{5.3cm} | X |} \hline
\rowcolor{gray!33} Object Field & Value Type & Description \\ \hline
"error\allowbreak{}Message" & \pref{String} &  \\ \hline
"error\allowbreak{}Code" & \pref{Number} &  \\ \hline
"exception\allowbreak{}Type" & \pref{ExceptionType} &  \\ \hline
"origin" & \pref{String} &  \\ \hline
\end{tabularx}



\msubsection{struct}{MetadataRequirement}

A JSON object of metadata requirements: each key is a metadata path (dot notation), each value the required value or an operator object such as \texttt{\{"op": "LESS\_\allowbreak{}THAN",\allowbreak{} "value": 10\}}.



\subsection{Primitives}
\label{sec:model:primitives}

As all messages are encoded using the JSON format \cite{rfc8259}, the following primitive constructs, part of that standard, become available.
Note that the official standard is defined in terms of parsing rules, while this list only concerns syntactic information.
Furthermore, the \pref{Object} and \pref{Array} types are given optional generic type parameters, which are used in this document to signify when pair values or elements are expected to conform to certain types.

\noindent\begin{tabularx}{\textwidth}{| p{4.6cm} | X |} \hline
\rowcolor{gray!33} Type & Description \\ \hline
\pdef{Array} & A JSON Array \cite{rfc8259}; \pref{Array}$<$A$>$ has elements of type A. \\ \hline
\pdef{AuthorizationLevel} & One of \texttt{MGMT}, \texttt{PROVIDER} (JSON String). \\ \hline
\pdef{AuthorizationPolicyType} & One of \texttt{ALL}, \texttt{WHITELIST}, \texttt{BLACKLIST}, \texttt{SYS\_\allowbreak{}METADATA} (JSON String). \\ \hline
\pdef{AuthorizationTargetType} & One of \texttt{SERVICE\_\allowbreak{}DEF}, \texttt{EVENT\_\allowbreak{}TYPE} (JSON String). \\ \hline
\pdef{Boolean} & JSON Boolean, \texttt{true} or \texttt{false}. \\ \hline
\pdef{DateTime} & JSON String with an ISO 8601 UTC timestamp, e.g. \texttt{2026-10-22T12:00:00Z}. \\ \hline
\pdef{ExceptionType} & One of \texttt{ARROWHEAD}, \texttt{INVALID\_\allowbreak{}PARAMETER}, \texttt{AUTH}, \texttt{FORBIDDEN}, \texttt{DATA\_\allowbreak{}NOT\_\allowbreak{}FOUND}, \texttt{TIMEOUT}, \texttt{LOCKED}, \texttt{INTERNAL\_\allowbreak{}SERVER\_\allowbreak{}ERROR}, \texttt{EXTERNAL\_\allowbreak{}SERVER\_\allowbreak{}ERROR} (JSON String). \\ \hline
\pdef{Number} & JSON Number, an integer. \\ \hline
\pdef{Object} & JSON Object; \pref{Object}$<$A$>$ has string keys and values of type A. \\ \hline
\pdef{String} & JSON String. \\ \hline
\end{tabularx}

\newpage

\bibliographystyle{IEEEtran}
\bibliography{bibliography}

\newpage

\section{Revision History}
\subsection{Amendments}

\noindent\begin{tabularx}{\textwidth}{| p{1cm} | p{3cm} | p{2cm} | X | p{4cm} |} \hline
\rowcolor{gray!33} No. & Date & Version & Subject of Amendments & Author \\ \hline

1 & 2026-10-22 & \arrowversion & Document created from the Arrowhead LaTeX templates for the 5.2.2 Java Spring implementation, including the local extensions of this workspace. & Jerker Delsing \\ \hline
\end{tabularx}

\subsection{Quality Assurance}

\noindent\begin{tabularx}{\textwidth}{| p{1cm} | p{3cm} | p{2cm} | X |} \hline
\rowcolor{gray!33} No. & Date & Version & Approved by \\ \hline

1 & & \arrowversion  &  \\ \hline

\end{tabularx}

\end{document}
